% ---------------------------------------------------------------------------
% Encryption-complexity comparison of the seven schemes under a FIXED client
% set (N, K treated as constants), expressed in PACK units:
%     cost = (# encrypted packs) x (cost per pack).
%
% ECPFL carries two complexity columns (no-dropout / dropout) because it is
% the only scheme whose encryption workload depends on dropout:
%   no-dropout : 1 vote ct (Eq. 9)  + k pack cts (Eq. 11)
%   dropout    : the above          + k compensatory cts (Eq. 14, one per
%                surviving client per consensus pack)
% per ckks_fed/ecpfl_crypto.py and main_ecpfl.py (vote/pack/comp loops).
%
% Requires : booktabs, amsmath
% Usage    : % ---------------------------------------------------------------------------
% Encryption-complexity comparison of the seven schemes under a FIXED client
% set (N, K treated as constants), expressed in PACK units:
%     cost = (# encrypted packs) x (cost per pack).
%
% ECPFL carries two complexity columns (no-dropout / dropout) because it is
% the only scheme whose encryption workload depends on dropout:
%   no-dropout : 1 vote ct (Eq. 9)  + k pack cts (Eq. 11)
%   dropout    : the above          + k compensatory cts (Eq. 14, one per
%                surviving client per consensus pack)
% per ckks_fed/ecpfl_crypto.py and main_ecpfl.py (vote/pack/comp loops).
%
% Requires : booktabs, amsmath
% Usage    : % ---------------------------------------------------------------------------
% Encryption-complexity comparison of the seven schemes under a FIXED client
% set (N, K treated as constants), expressed in PACK units:
%     cost = (# encrypted packs) x (cost per pack).
%
% ECPFL carries two complexity columns (no-dropout / dropout) because it is
% the only scheme whose encryption workload depends on dropout:
%   no-dropout : 1 vote ct (Eq. 9)  + k pack cts (Eq. 11)
%   dropout    : the above          + k compensatory cts (Eq. 14, one per
%                surviving client per consensus pack)
% per ckks_fed/ecpfl_crypto.py and main_ecpfl.py (vote/pack/comp loops).
%
% Requires : booktabs, amsmath
% Usage    : % ---------------------------------------------------------------------------
% Encryption-complexity comparison of the seven schemes under a FIXED client
% set (N, K treated as constants), expressed in PACK units:
%     cost = (# encrypted packs) x (cost per pack).
%
% ECPFL carries two complexity columns (no-dropout / dropout) because it is
% the only scheme whose encryption workload depends on dropout:
%   no-dropout : 1 vote ct (Eq. 9)  + k pack cts (Eq. 11)
%   dropout    : the above          + k compensatory cts (Eq. 14, one per
%                surviving client per consensus pack)
% per ckks_fed/ecpfl_crypto.py and main_ecpfl.py (vote/pack/comp loops).
%
% Requires : booktabs, amsmath
% Usage    : \input{encryption_complexity_table.tex}
% Symbols  : follow tab:comp-notation in computational_complexity.tex
%            (d, n, J, s, k, L, N, K, O_ntt)
% ---------------------------------------------------------------------------
\begin{table}[t]
\centering
\caption{Client-side encryption complexity per round under a fixed client
set, expressed in \emph{pack units}: cost $=$ (number of encrypted packs)
$\times$ (cost per pack). $J$ is the number of packs covering the full
model under each scheme's own packing, $s\in(0,1]$ the FedPHE sparsity
ratio, $k\le J$ the number of ECPFL consensus packs, and $L$ the number of
RNS moduli of multi-key CKKS; $\kappa$ is the Paillier modulus bit length
and $m=\big\lfloor(\kappa-1)/\lceil\log_2(nQ)\rceil\big\rfloor$ the number
of sample points packed into one Paillier ciphertext (SAMK). The four
baselines FedAvg+CKKS, FedPHE, MASER and ESAFL provide no dropout
tolerance ($\times$ in the last column): their per-client encryption cost
is only defined for a dropout-free round; SAMK's cost is dropout-invariant
(survivors finish the round, the dropped client forfeits it by design);
SecFed and ECPFL complete dropout rounds---SecFed via TEE-generated
supplementary ciphertexts (no extra client encryption), ECPFL via
compensatory encryptions paid by the surviving clients.}
\label{tab:encryption-complexity}
\small
\renewcommand{\arraystretch}{1.22}
\begin{tabular}{@{}lcccc@{}}
\toprule
& & & \multicolumn{2}{c}{\textbf{Encryption complexity}} \\
\cmidrule(lr){4-5}
\textbf{Scheme} & \textbf{Packs enc.} & \textbf{Cost / pack} &
\textbf{No dropout} & \textbf{Dropout} \\
\midrule
FedAvg+CKKS  & $J$    & $\mathcal{O}(n\log n)$    &
$J\cdot\mathcal{O}(n\log n)$ & $\times$ \\[3pt]
FedPHE       & $sJ$   & $\mathcal{O}(n\log n)$    &
$sJ\cdot\mathcal{O}(n\log n)$ & $\times$ \\[3pt]
MASER        & $J$    & $\mathcal{O}(n\log n+nL)$ &
$J\cdot\mathcal{O}(n\log n+nL)$ & $\times$ \\[3pt]
ESAFL        & $J$    & $\mathcal{O}(n\log n)$    &
$J\cdot\mathcal{O}(n\log n)$ & $\times$ \\[3pt]
SecFed       & $J$    & $\mathcal{O}(n\log n)$    &
$J\cdot\mathcal{O}(n\log n)$ & (same) \\[3pt]
SAMK         & \multicolumn{2}{c}{$J$ chunks:
$1$ BFV enc.\ $+$ $\lceil\frac{n-1}{m}\rceil{+}N$ Paillier enc.\ each} &
$J\cdot\mathcal{O}\!\left(n^{2}+\big(\lceil\tfrac{n-1}{m}\rceil{+}N\big)\kappa^{3}\right)$
& (same) \\[3pt]
\midrule
ECPFL (ours) & $1+k\;(+k)$ & $\mathcal{O}(n\log n)$ &
$(k{+}1)\cdot\mathcal{O}(n\log n)+\mathcal{O}(Nn)$ &
$(2k{+}1)\cdot\mathcal{O}(n\log n)+\mathcal{O}(Nn)$ \\
\bottomrule
\end{tabular}

\vspace{4pt}
\begin{minipage}{\linewidth}
\footnotesize
\emph{Notes.}
(i) \emph{Packs enc.} for ECPFL $=$ one vote ciphertext (the whole $J$-bit
bitmap packed into the low coefficients of a single ring element) $+$ $k$
consensus-pack ciphertexts; the parenthesized $(+k)$ are compensatory
ciphertexts, executed by each \emph{surviving} client per consensus pack on
dropout rounds only.
(ii) $\times$ marks the absence of dropout tolerance, with
scheme-specific failure modes: with FedAvg+CKKS the dropped updates are
simply lost and the round still completes for the survivors; with FedPHE
the dropped client's obfuscating factor $\Delta_i$ (dispatched by the KDC)
remains uncancelled in the aggregate, so deobfuscation fails for
\emph{every} client and the round is invalid; with MASER both the
majority-vote mask consensus and the xMK-CKKS decryption
($\mathrm{PartialDec}$ by \emph{every} participant $+\ \mathrm{Merge}$)
require the full participant set; with ESAFL the summed key
$S=\sum_i s_i$ requires all contributions---in the last three cases the
round cannot complete at all, and in none of them is a meaningful
``dropout encryption complexity'' defined; ECPFL is the only scheme that
pays additional encryptions in order to finish the round.
(iii) The additive $\mathcal{O}(Nn)$ term is the share-sum key material
$\tilde{s}_{i,j}=s_i+\lambda_i\sum_{k\notin\mathcal{S}}s_{k,i}
+\lambda_i b_j z_i$ (Eq.~11), paid once per round and cached across the $k$
packs; with $N$ fixed it is a constant multiple of $n$.
(iv) Pack units keep different packing densities comparable:
$J=\lceil 2d/n\rceil$ for slot-packed CKKS, $\lceil d/2n\rceil$ for ESAFL's
EP packing, $\lceil d/n\rceil$ for ECPFL's, SecFed's and SAMK's
coefficient packing (SAMK uses its own ring degree, $n=1024$ in the
reproduction); in
parameter units every $J$-term equals $\mathcal{O}(d\log n)$, and MASER
additionally pays the RNS term $nL$ per pack.
(v) SecFed (TDSC 2024, EMK-BFV), per pack (Eq.~12):
$\mathsf{ct}[0]=[\mathrm{pk}[0]\cdot\mu+e_1+\Delta\cdot\mathrm{pt}]_q$ and
$\mathsf{ct}[1]=[\mathrm{pk}[1]\cdot\mu+e_2]_q$---two ring multiplications
per pack, the same two-multiplication structure as ECPFL. Moreover,
$\mathrm{pk}[0]\cdot\mu$ and $\mathrm{pk}[1]\cdot\mu$ are \emph{per-user
constants} (the public key and the system parameter $\mu$ are fixed), so if
cached across packs the marginal per-pack cost drops to $\mathcal{O}(n)$
and the per-round total to $\mathcal{O}(n\log n+d)$, a $\log$-factor below
$J\cdot\mathcal{O}(n\log n)$; the table reports the uncached count for
uniformity. Its dropout column reads ``(same)'': the offline protection
mechanism (Algorithm~8, \emph{supEnc}) runs in the TEE, which generates
supplementary ciphertexts for the offline users under their public keys
(tolerating up to $N-2$ dropouts) at no extra client-side encryption.
(vi) SAMK (TDSC 2025), per chunk: one BFV encryption
($\mathcal{O}(n\log n)$, one ring multiplication), one ring product
$f_i=\mathrm{sk}_i\cdot c_{i,1}$ ($\mathcal{O}(n\log n)$), Horner
evaluation of $f_i$ at the $n$ interpolation points ($\mathcal{O}(n^{2})$),
and $\lceil(n-1)/m\rceil+N$ Paillier encryptions: the first $n-1$ points
packed under the server's public key ($m$ points per ciphertext), and the
last point separately under \emph{every} client's public key for access
control. One Paillier encryption is a $\kappa$-bit modular exponentiation,
$\mathcal{O}(\kappa^{3})$ bit operations---a different unit from the
word-level ring operations, so SAMK's dominant term is not directly
comparable in magnitude but is orders of magnitude heavier in practice.
Its dropout column reads ``(same)'': survivors complete the round at
unchanged cost, while the dropped client forfeits the round
($\mathsf{SAMKAccessError}$)---access control by design, not compensation.
(vii) Each ECPFL pack encryption performs two ring multiplications
($a_j\tilde{s}_{i,j}$ and $b_j z_i$) against one for the baselines, so
ECPFL is cheaper than full-model encryption only while $2k+1<J$, i.e.\
$\theta<\tfrac12$ under rank consensus $k=\lceil\theta J\rceil$.
\end{minipage}
\end{table}

% Symbols  : follow tab:comp-notation in computational_complexity.tex
%            (d, n, J, s, k, L, N, K, O_ntt)
% ---------------------------------------------------------------------------
\begin{table}[t]
\centering
\caption{Client-side encryption complexity per round under a fixed client
set, expressed in \emph{pack units}: cost $=$ (number of encrypted packs)
$\times$ (cost per pack). $J$ is the number of packs covering the full
model under each scheme's own packing, $s\in(0,1]$ the FedPHE sparsity
ratio, $k\le J$ the number of ECPFL consensus packs, and $L$ the number of
RNS moduli of multi-key CKKS; $\kappa$ is the Paillier modulus bit length
and $m=\big\lfloor(\kappa-1)/\lceil\log_2(nQ)\rceil\big\rfloor$ the number
of sample points packed into one Paillier ciphertext (SAMK). The four
baselines FedAvg+CKKS, FedPHE, MASER and ESAFL provide no dropout
tolerance ($\times$ in the last column): their per-client encryption cost
is only defined for a dropout-free round; SAMK's cost is dropout-invariant
(survivors finish the round, the dropped client forfeits it by design);
SecFed and ECPFL complete dropout rounds---SecFed via TEE-generated
supplementary ciphertexts (no extra client encryption), ECPFL via
compensatory encryptions paid by the surviving clients.}
\label{tab:encryption-complexity}
\small
\renewcommand{\arraystretch}{1.22}
\begin{tabular}{@{}lcccc@{}}
\toprule
& & & \multicolumn{2}{c}{\textbf{Encryption complexity}} \\
\cmidrule(lr){4-5}
\textbf{Scheme} & \textbf{Packs enc.} & \textbf{Cost / pack} &
\textbf{No dropout} & \textbf{Dropout} \\
\midrule
FedAvg+CKKS  & $J$    & $\mathcal{O}(n\log n)$    &
$J\cdot\mathcal{O}(n\log n)$ & $\times$ \\[3pt]
FedPHE       & $sJ$   & $\mathcal{O}(n\log n)$    &
$sJ\cdot\mathcal{O}(n\log n)$ & $\times$ \\[3pt]
MASER        & $J$    & $\mathcal{O}(n\log n+nL)$ &
$J\cdot\mathcal{O}(n\log n+nL)$ & $\times$ \\[3pt]
ESAFL        & $J$    & $\mathcal{O}(n\log n)$    &
$J\cdot\mathcal{O}(n\log n)$ & $\times$ \\[3pt]
SecFed       & $J$    & $\mathcal{O}(n\log n)$    &
$J\cdot\mathcal{O}(n\log n)$ & (same) \\[3pt]
SAMK         & \multicolumn{2}{c}{$J$ chunks:
$1$ BFV enc.\ $+$ $\lceil\frac{n-1}{m}\rceil{+}N$ Paillier enc.\ each} &
$J\cdot\mathcal{O}\!\left(n^{2}+\big(\lceil\tfrac{n-1}{m}\rceil{+}N\big)\kappa^{3}\right)$
& (same) \\[3pt]
\midrule
ECPFL (ours) & $1+k\;(+k)$ & $\mathcal{O}(n\log n)$ &
$(k{+}1)\cdot\mathcal{O}(n\log n)+\mathcal{O}(Nn)$ &
$(2k{+}1)\cdot\mathcal{O}(n\log n)+\mathcal{O}(Nn)$ \\
\bottomrule
\end{tabular}

\vspace{4pt}
\begin{minipage}{\linewidth}
\footnotesize
\emph{Notes.}
(i) \emph{Packs enc.} for ECPFL $=$ one vote ciphertext (the whole $J$-bit
bitmap packed into the low coefficients of a single ring element) $+$ $k$
consensus-pack ciphertexts; the parenthesized $(+k)$ are compensatory
ciphertexts, executed by each \emph{surviving} client per consensus pack on
dropout rounds only.
(ii) $\times$ marks the absence of dropout tolerance, with
scheme-specific failure modes: with FedAvg+CKKS the dropped updates are
simply lost and the round still completes for the survivors; with FedPHE
the dropped client's obfuscating factor $\Delta_i$ (dispatched by the KDC)
remains uncancelled in the aggregate, so deobfuscation fails for
\emph{every} client and the round is invalid; with MASER both the
majority-vote mask consensus and the xMK-CKKS decryption
($\mathrm{PartialDec}$ by \emph{every} participant $+\ \mathrm{Merge}$)
require the full participant set; with ESAFL the summed key
$S=\sum_i s_i$ requires all contributions---in the last three cases the
round cannot complete at all, and in none of them is a meaningful
``dropout encryption complexity'' defined; ECPFL is the only scheme that
pays additional encryptions in order to finish the round.
(iii) The additive $\mathcal{O}(Nn)$ term is the share-sum key material
$\tilde{s}_{i,j}=s_i+\lambda_i\sum_{k\notin\mathcal{S}}s_{k,i}
+\lambda_i b_j z_i$ (Eq.~11), paid once per round and cached across the $k$
packs; with $N$ fixed it is a constant multiple of $n$.
(iv) Pack units keep different packing densities comparable:
$J=\lceil 2d/n\rceil$ for slot-packed CKKS, $\lceil d/2n\rceil$ for ESAFL's
EP packing, $\lceil d/n\rceil$ for ECPFL's, SecFed's and SAMK's
coefficient packing (SAMK uses its own ring degree, $n=1024$ in the
reproduction); in
parameter units every $J$-term equals $\mathcal{O}(d\log n)$, and MASER
additionally pays the RNS term $nL$ per pack.
(v) SecFed (TDSC 2024, EMK-BFV), per pack (Eq.~12):
$\mathsf{ct}[0]=[\mathrm{pk}[0]\cdot\mu+e_1+\Delta\cdot\mathrm{pt}]_q$ and
$\mathsf{ct}[1]=[\mathrm{pk}[1]\cdot\mu+e_2]_q$---two ring multiplications
per pack, the same two-multiplication structure as ECPFL. Moreover,
$\mathrm{pk}[0]\cdot\mu$ and $\mathrm{pk}[1]\cdot\mu$ are \emph{per-user
constants} (the public key and the system parameter $\mu$ are fixed), so if
cached across packs the marginal per-pack cost drops to $\mathcal{O}(n)$
and the per-round total to $\mathcal{O}(n\log n+d)$, a $\log$-factor below
$J\cdot\mathcal{O}(n\log n)$; the table reports the uncached count for
uniformity. Its dropout column reads ``(same)'': the offline protection
mechanism (Algorithm~8, \emph{supEnc}) runs in the TEE, which generates
supplementary ciphertexts for the offline users under their public keys
(tolerating up to $N-2$ dropouts) at no extra client-side encryption.
(vi) SAMK (TDSC 2025), per chunk: one BFV encryption
($\mathcal{O}(n\log n)$, one ring multiplication), one ring product
$f_i=\mathrm{sk}_i\cdot c_{i,1}$ ($\mathcal{O}(n\log n)$), Horner
evaluation of $f_i$ at the $n$ interpolation points ($\mathcal{O}(n^{2})$),
and $\lceil(n-1)/m\rceil+N$ Paillier encryptions: the first $n-1$ points
packed under the server's public key ($m$ points per ciphertext), and the
last point separately under \emph{every} client's public key for access
control. One Paillier encryption is a $\kappa$-bit modular exponentiation,
$\mathcal{O}(\kappa^{3})$ bit operations---a different unit from the
word-level ring operations, so SAMK's dominant term is not directly
comparable in magnitude but is orders of magnitude heavier in practice.
Its dropout column reads ``(same)'': survivors complete the round at
unchanged cost, while the dropped client forfeits the round
($\mathsf{SAMKAccessError}$)---access control by design, not compensation.
(vii) Each ECPFL pack encryption performs two ring multiplications
($a_j\tilde{s}_{i,j}$ and $b_j z_i$) against one for the baselines, so
ECPFL is cheaper than full-model encryption only while $2k+1<J$, i.e.\
$\theta<\tfrac12$ under rank consensus $k=\lceil\theta J\rceil$.
\end{minipage}
\end{table}

% Symbols  : follow tab:comp-notation in computational_complexity.tex
%            (d, n, J, s, k, L, N, K, O_ntt)
% ---------------------------------------------------------------------------
\begin{table}[t]
\centering
\caption{Client-side encryption complexity per round under a fixed client
set, expressed in \emph{pack units}: cost $=$ (number of encrypted packs)
$\times$ (cost per pack). $J$ is the number of packs covering the full
model under each scheme's own packing, $s\in(0,1]$ the FedPHE sparsity
ratio, $k\le J$ the number of ECPFL consensus packs, and $L$ the number of
RNS moduli of multi-key CKKS; $\kappa$ is the Paillier modulus bit length
and $m=\big\lfloor(\kappa-1)/\lceil\log_2(nQ)\rceil\big\rfloor$ the number
of sample points packed into one Paillier ciphertext (SAMK). The four
baselines FedAvg+CKKS, FedPHE, MASER and ESAFL provide no dropout
tolerance ($\times$ in the last column): their per-client encryption cost
is only defined for a dropout-free round; SAMK's cost is dropout-invariant
(survivors finish the round, the dropped client forfeits it by design);
SecFed and ECPFL complete dropout rounds---SecFed via TEE-generated
supplementary ciphertexts (no extra client encryption), ECPFL via
compensatory encryptions paid by the surviving clients.}
\label{tab:encryption-complexity}
\small
\renewcommand{\arraystretch}{1.22}
\begin{tabular}{@{}lcccc@{}}
\toprule
& & & \multicolumn{2}{c}{\textbf{Encryption complexity}} \\
\cmidrule(lr){4-5}
\textbf{Scheme} & \textbf{Packs enc.} & \textbf{Cost / pack} &
\textbf{No dropout} & \textbf{Dropout} \\
\midrule
FedAvg+CKKS  & $J$    & $\mathcal{O}(n\log n)$    &
$J\cdot\mathcal{O}(n\log n)$ & $\times$ \\[3pt]
FedPHE       & $sJ$   & $\mathcal{O}(n\log n)$    &
$sJ\cdot\mathcal{O}(n\log n)$ & $\times$ \\[3pt]
MASER        & $J$    & $\mathcal{O}(n\log n+nL)$ &
$J\cdot\mathcal{O}(n\log n+nL)$ & $\times$ \\[3pt]
ESAFL        & $J$    & $\mathcal{O}(n\log n)$    &
$J\cdot\mathcal{O}(n\log n)$ & $\times$ \\[3pt]
SecFed       & $J$    & $\mathcal{O}(n\log n)$    &
$J\cdot\mathcal{O}(n\log n)$ & (same) \\[3pt]
SAMK         & \multicolumn{2}{c}{$J$ chunks:
$1$ BFV enc.\ $+$ $\lceil\frac{n-1}{m}\rceil{+}N$ Paillier enc.\ each} &
$J\cdot\mathcal{O}\!\left(n^{2}+\big(\lceil\tfrac{n-1}{m}\rceil{+}N\big)\kappa^{3}\right)$
& (same) \\[3pt]
\midrule
ECPFL (ours) & $1+k\;(+k)$ & $\mathcal{O}(n\log n)$ &
$(k{+}1)\cdot\mathcal{O}(n\log n)+\mathcal{O}(Nn)$ &
$(2k{+}1)\cdot\mathcal{O}(n\log n)+\mathcal{O}(Nn)$ \\
\bottomrule
\end{tabular}

\vspace{4pt}
\begin{minipage}{\linewidth}
\footnotesize
\emph{Notes.}
(i) \emph{Packs enc.} for ECPFL $=$ one vote ciphertext (the whole $J$-bit
bitmap packed into the low coefficients of a single ring element) $+$ $k$
consensus-pack ciphertexts; the parenthesized $(+k)$ are compensatory
ciphertexts, executed by each \emph{surviving} client per consensus pack on
dropout rounds only.
(ii) $\times$ marks the absence of dropout tolerance, with
scheme-specific failure modes: with FedAvg+CKKS the dropped updates are
simply lost and the round still completes for the survivors; with FedPHE
the dropped client's obfuscating factor $\Delta_i$ (dispatched by the KDC)
remains uncancelled in the aggregate, so deobfuscation fails for
\emph{every} client and the round is invalid; with MASER both the
majority-vote mask consensus and the xMK-CKKS decryption
($\mathrm{PartialDec}$ by \emph{every} participant $+\ \mathrm{Merge}$)
require the full participant set; with ESAFL the summed key
$S=\sum_i s_i$ requires all contributions---in the last three cases the
round cannot complete at all, and in none of them is a meaningful
``dropout encryption complexity'' defined; ECPFL is the only scheme that
pays additional encryptions in order to finish the round.
(iii) The additive $\mathcal{O}(Nn)$ term is the share-sum key material
$\tilde{s}_{i,j}=s_i+\lambda_i\sum_{k\notin\mathcal{S}}s_{k,i}
+\lambda_i b_j z_i$ (Eq.~11), paid once per round and cached across the $k$
packs; with $N$ fixed it is a constant multiple of $n$.
(iv) Pack units keep different packing densities comparable:
$J=\lceil 2d/n\rceil$ for slot-packed CKKS, $\lceil d/2n\rceil$ for ESAFL's
EP packing, $\lceil d/n\rceil$ for ECPFL's, SecFed's and SAMK's
coefficient packing (SAMK uses its own ring degree, $n=1024$ in the
reproduction); in
parameter units every $J$-term equals $\mathcal{O}(d\log n)$, and MASER
additionally pays the RNS term $nL$ per pack.
(v) SecFed (TDSC 2024, EMK-BFV), per pack (Eq.~12):
$\mathsf{ct}[0]=[\mathrm{pk}[0]\cdot\mu+e_1+\Delta\cdot\mathrm{pt}]_q$ and
$\mathsf{ct}[1]=[\mathrm{pk}[1]\cdot\mu+e_2]_q$---two ring multiplications
per pack, the same two-multiplication structure as ECPFL. Moreover,
$\mathrm{pk}[0]\cdot\mu$ and $\mathrm{pk}[1]\cdot\mu$ are \emph{per-user
constants} (the public key and the system parameter $\mu$ are fixed), so if
cached across packs the marginal per-pack cost drops to $\mathcal{O}(n)$
and the per-round total to $\mathcal{O}(n\log n+d)$, a $\log$-factor below
$J\cdot\mathcal{O}(n\log n)$; the table reports the uncached count for
uniformity. Its dropout column reads ``(same)'': the offline protection
mechanism (Algorithm~8, \emph{supEnc}) runs in the TEE, which generates
supplementary ciphertexts for the offline users under their public keys
(tolerating up to $N-2$ dropouts) at no extra client-side encryption.
(vi) SAMK (TDSC 2025), per chunk: one BFV encryption
($\mathcal{O}(n\log n)$, one ring multiplication), one ring product
$f_i=\mathrm{sk}_i\cdot c_{i,1}$ ($\mathcal{O}(n\log n)$), Horner
evaluation of $f_i$ at the $n$ interpolation points ($\mathcal{O}(n^{2})$),
and $\lceil(n-1)/m\rceil+N$ Paillier encryptions: the first $n-1$ points
packed under the server's public key ($m$ points per ciphertext), and the
last point separately under \emph{every} client's public key for access
control. One Paillier encryption is a $\kappa$-bit modular exponentiation,
$\mathcal{O}(\kappa^{3})$ bit operations---a different unit from the
word-level ring operations, so SAMK's dominant term is not directly
comparable in magnitude but is orders of magnitude heavier in practice.
Its dropout column reads ``(same)'': survivors complete the round at
unchanged cost, while the dropped client forfeits the round
($\mathsf{SAMKAccessError}$)---access control by design, not compensation.
(vii) Each ECPFL pack encryption performs two ring multiplications
($a_j\tilde{s}_{i,j}$ and $b_j z_i$) against one for the baselines, so
ECPFL is cheaper than full-model encryption only while $2k+1<J$, i.e.\
$\theta<\tfrac12$ under rank consensus $k=\lceil\theta J\rceil$.
\end{minipage}
\end{table}

% Symbols  : follow tab:comp-notation in computational_complexity.tex
%            (d, n, J, s, k, L, N, K, O_ntt)
% ---------------------------------------------------------------------------
\begin{table}[t]
\centering
\caption{Client-side encryption complexity per round under a fixed client
set, expressed in \emph{pack units}: cost $=$ (number of encrypted packs)
$\times$ (cost per pack). $J$ is the number of packs covering the full
model under each scheme's own packing, $s\in(0,1]$ the FedPHE sparsity
ratio, $k\le J$ the number of ECPFL consensus packs, and $L$ the number of
RNS moduli of multi-key CKKS; $\kappa$ is the Paillier modulus bit length
and $m=\big\lfloor(\kappa-1)/\lceil\log_2(nQ)\rceil\big\rfloor$ the number
of sample points packed into one Paillier ciphertext (SAMK). The four
baselines FedAvg+CKKS, FedPHE, MASER and ESAFL provide no dropout
tolerance ($\times$ in the last column): their per-client encryption cost
is only defined for a dropout-free round; SAMK's cost is dropout-invariant
(survivors finish the round, the dropped client forfeits it by design);
SecFed and ECPFL complete dropout rounds---SecFed via TEE-generated
supplementary ciphertexts (no extra client encryption), ECPFL via
compensatory encryptions paid by the surviving clients.}
\label{tab:encryption-complexity}
\small
\renewcommand{\arraystretch}{1.22}
\begin{tabular}{@{}lcccc@{}}
\toprule
& & & \multicolumn{2}{c}{\textbf{Encryption complexity}} \\
\cmidrule(lr){4-5}
\textbf{Scheme} & \textbf{Packs enc.} & \textbf{Cost / pack} &
\textbf{No dropout} & \textbf{Dropout} \\
\midrule
FedAvg+CKKS  & $J$    & $\mathcal{O}(n\log n)$    &
$J\cdot\mathcal{O}(n\log n)$ & $\times$ \\[3pt]
FedPHE       & $sJ$   & $\mathcal{O}(n\log n)$    &
$sJ\cdot\mathcal{O}(n\log n)$ & $\times$ \\[3pt]
MASER        & $J$    & $\mathcal{O}(n\log n+nL)$ &
$J\cdot\mathcal{O}(n\log n+nL)$ & $\times$ \\[3pt]
ESAFL        & $J$    & $\mathcal{O}(n\log n)$    &
$J\cdot\mathcal{O}(n\log n)$ & $\times$ \\[3pt]
SecFed       & $J$    & $\mathcal{O}(n\log n)$    &
$J\cdot\mathcal{O}(n\log n)$ & (same) \\[3pt]
SAMK         & \multicolumn{2}{c}{$J$ chunks:
$1$ BFV enc.\ $+$ $\lceil\frac{n-1}{m}\rceil{+}N$ Paillier enc.\ each} &
$J\cdot\mathcal{O}\!\left(n^{2}+\big(\lceil\tfrac{n-1}{m}\rceil{+}N\big)\kappa^{3}\right)$
& (same) \\[3pt]
\midrule
ECPFL (ours) & $1+k\;(+k)$ & $\mathcal{O}(n\log n)$ &
$(k{+}1)\cdot\mathcal{O}(n\log n)+\mathcal{O}(Nn)$ &
$(2k{+}1)\cdot\mathcal{O}(n\log n)+\mathcal{O}(Nn)$ \\
\bottomrule
\end{tabular}

\vspace{4pt}
\begin{minipage}{\linewidth}
\footnotesize
\emph{Notes.}
(i) \emph{Packs enc.} for ECPFL $=$ one vote ciphertext (the whole $J$-bit
bitmap packed into the low coefficients of a single ring element) $+$ $k$
consensus-pack ciphertexts; the parenthesized $(+k)$ are compensatory
ciphertexts, executed by each \emph{surviving} client per consensus pack on
dropout rounds only.
(ii) $\times$ marks the absence of dropout tolerance, with
scheme-specific failure modes: with FedAvg+CKKS the dropped updates are
simply lost and the round still completes for the survivors; with FedPHE
the dropped client's obfuscating factor $\Delta_i$ (dispatched by the KDC)
remains uncancelled in the aggregate, so deobfuscation fails for
\emph{every} client and the round is invalid; with MASER both the
majority-vote mask consensus and the xMK-CKKS decryption
($\mathrm{PartialDec}$ by \emph{every} participant $+\ \mathrm{Merge}$)
require the full participant set; with ESAFL the summed key
$S=\sum_i s_i$ requires all contributions---in the last three cases the
round cannot complete at all, and in none of them is a meaningful
``dropout encryption complexity'' defined; ECPFL is the only scheme that
pays additional encryptions in order to finish the round.
(iii) The additive $\mathcal{O}(Nn)$ term is the share-sum key material
$\tilde{s}_{i,j}=s_i+\lambda_i\sum_{k\notin\mathcal{S}}s_{k,i}
+\lambda_i b_j z_i$ (Eq.~11), paid once per round and cached across the $k$
packs; with $N$ fixed it is a constant multiple of $n$.
(iv) Pack units keep different packing densities comparable:
$J=\lceil 2d/n\rceil$ for slot-packed CKKS, $\lceil d/2n\rceil$ for ESAFL's
EP packing, $\lceil d/n\rceil$ for ECPFL's, SecFed's and SAMK's
coefficient packing (SAMK uses its own ring degree, $n=1024$ in the
reproduction); in
parameter units every $J$-term equals $\mathcal{O}(d\log n)$, and MASER
additionally pays the RNS term $nL$ per pack.
(v) SecFed (TDSC 2024, EMK-BFV), per pack (Eq.~12):
$\mathsf{ct}[0]=[\mathrm{pk}[0]\cdot\mu+e_1+\Delta\cdot\mathrm{pt}]_q$ and
$\mathsf{ct}[1]=[\mathrm{pk}[1]\cdot\mu+e_2]_q$---two ring multiplications
per pack, the same two-multiplication structure as ECPFL. Moreover,
$\mathrm{pk}[0]\cdot\mu$ and $\mathrm{pk}[1]\cdot\mu$ are \emph{per-user
constants} (the public key and the system parameter $\mu$ are fixed), so if
cached across packs the marginal per-pack cost drops to $\mathcal{O}(n)$
and the per-round total to $\mathcal{O}(n\log n+d)$, a $\log$-factor below
$J\cdot\mathcal{O}(n\log n)$; the table reports the uncached count for
uniformity. Its dropout column reads ``(same)'': the offline protection
mechanism (Algorithm~8, \emph{supEnc}) runs in the TEE, which generates
supplementary ciphertexts for the offline users under their public keys
(tolerating up to $N-2$ dropouts) at no extra client-side encryption.
(vi) SAMK (TDSC 2025), per chunk: one BFV encryption
($\mathcal{O}(n\log n)$, one ring multiplication), one ring product
$f_i=\mathrm{sk}_i\cdot c_{i,1}$ ($\mathcal{O}(n\log n)$), Horner
evaluation of $f_i$ at the $n$ interpolation points ($\mathcal{O}(n^{2})$),
and $\lceil(n-1)/m\rceil+N$ Paillier encryptions: the first $n-1$ points
packed under the server's public key ($m$ points per ciphertext), and the
last point separately under \emph{every} client's public key for access
control. One Paillier encryption is a $\kappa$-bit modular exponentiation,
$\mathcal{O}(\kappa^{3})$ bit operations---a different unit from the
word-level ring operations, so SAMK's dominant term is not directly
comparable in magnitude but is orders of magnitude heavier in practice.
Its dropout column reads ``(same)'': survivors complete the round at
unchanged cost, while the dropped client forfeits the round
($\mathsf{SAMKAccessError}$)---access control by design, not compensation.
(vii) Each ECPFL pack encryption performs two ring multiplications
($a_j\tilde{s}_{i,j}$ and $b_j z_i$) against one for the baselines, so
ECPFL is cheaper than full-model encryption only while $2k+1<J$, i.e.\
$\theta<\tfrac12$ under rank consensus $k=\lceil\theta J\rceil$.
\end{minipage}
\end{table}
